\chapter{Number Fields and Their Integers}
\begin{center}\itshape
Is $\Z[\sqrt[3]2]$ the ring of all integers of $\Q(\sqrt[3]2)$? And $\Z[\sqrt[3]{10}]$ for $\Q(\sqrt[3]{10})$?\\
The first answer is yes and the second is no, and this chapter explains how one finds out.
\end{center}

In 1770 Leonhard Euler, by then nearly blind, dictated an algebra textbook to a servant, and in it he proved that
$x^3+y^3=z^3$ has no solution in positive integers. At the crucial step he needed to know when $a^2+3b^2$ is a cube, and
he simply wrote $a^2+3b^2=(a+b\sqrt{-3})(a-b\sqrt{-3})$ and declared that each factor had to be a cube too, ``as with
ordinary numbers''. The theorem was right, the step was not quite: in $\Z[\sqrt{-3}]$ the number $4=2\cdot2=(1+\sqrt{-3})
(1-\sqrt{-3})$ factors in two different ways. Euler had walked through a door without noticing that it was the wrong one.

The repair took a century, and its first lesson is in Tome~1's quadratic fields. The ring to work in is not $\Z[\sqrt{-3}]$
but the slightly larger $\Z[\frac{1+\sqrt{-3}}2]$, where factorisation is unique and Euler's argument goes through. Tome~1
found this ring by hand (Tome~1, Proposition~11.2): $\Z[\sqrt d]$ when $d\equiv2,3\pmod 4$, and the larger
$\Z[\frac{1+\sqrt d}2]$ when $d\equiv1\pmod4$. So the ``obvious'' ring $\Z[\alpha]$ generated by a defining number can be
too small. The right ring must be defined by a \emph{property}, not by a choice of generator. That property, being a root
of a monic polynomial with integer coefficients, is what Dedekind took as the definition of an algebraic integer in
1871, in a supplement he wrote for his late teacher Dirichlet's lectures.

This chapter sets up that world in any degree: number fields, their embeddings into $\C$, algebraic integers, norms,
traces and discriminants. The hero is a single theorem. The integers $\mathcal O_K$ of a number field of degree $n$ form
a \emph{lattice}, a free abelian group of rank $n$, and the volume of this lattice, the \emph{discriminant} $d_K$, turns
out to control almost everything that follows: which primes ramify (Chapter~2), how large the class group can be
(Chapter~3), the shape of the zeta function (Chapters~8--9). We finish with practical tools to compute $\mathcal O_K$,
enough to settle the two questions at the top of this page, and with a cubic field that Dedekind built in 1878 to show
that sometimes no single generator will do.

\begin{reminder}[from Tomes 0 and 2]
\emph{Minimal polynomials} (Tome~2, Proposition~8.5 and Theorem~8.6). If $\alpha\in\C$ is algebraic over $\Q$, the
polynomials of $\Q[X]$ vanishing at $\alpha$ are the multiples of a unique monic irreducible $m_\alpha$, and
$\Q(\alpha)=\Q[\alpha]\cong\Q[X]/(m_\alpha)$ has degree $\deg m_\alpha$, with basis $1,\alpha,\dots,\alpha^{d-1}$.
\emph{Separability} (Tome~2, Lemma~8.13). An irreducible polynomial over a field of characteristic $0$ has distinct
roots. \emph{Primitive elements} (Tome~2, Exercise~8.8). Every finite extension of $\Q$ inside $\C$ is $\Q(\theta)$
for some $\theta$. \emph{Rational roots} (Tome~0, Proposition~3.7). A rational root of a monic polynomial with
integer coefficients is an integer. \emph{Lattices} (Tome~2, Lemma~7.5 and Corollary~7.8). A subgroup of $\Z^n$ is free
of rank at most $n$; if $A\in M_n(\Z)$ has $\det A\ne0$, then $A\Z^n$ has index $|\det A|$ in $\Z^n$. \emph{Galois
theory} (Tome~2, Theorems~10.4 and~10.6). The splitting field $L\subseteq\C$ of a polynomial of $\Q[X]$ is Galois over
$\Q$, and the elements of $L$ fixed by every automorphism of $L$ are the rationals.
\end{reminder}

\section{Number fields and their embeddings}

\begin{definition}
A \emph{number field} is a subfield $K$ of $\C$ of finite degree $n=[K:\Q]$ over $\Q$. An \emph{embedding} of $K$ is
a ring homomorphism $\sigma:K\to\C$ (it fixes $\Q$, since it fixes $1$).
\end{definition}

Working inside $\C$ is a convenience, as in Tomes~1 and~2: every field of finite degree over $\Q$ is isomorphic to
such a $K$ (Tome~2, Theorem~8.9 and the fundamental theorem of algebra), and nothing below depends on the choice.

\begin{proposition}\label{prop:embeddings}
A number field $K$ of degree $n$ has exactly $n$ embeddings. If $K=\Q(\theta)$ and $\theta_1,\dots,\theta_n$ are the
roots of $m_\theta$, they are the maps $\sigma_i$ with $\sigma_i(\theta)=\theta_i$.
\end{proposition}
\begin{proof}
Write $K=\Q(\theta)$ (primitive element). The roots $\theta_i$ are distinct (separability), and there are $n$ of them
since $\deg m_\theta=[K:\Q]=n$. An embedding $\sigma$ satisfies $m_\theta(\sigma(\theta))=\sigma(m_\theta(\theta))=0$,
so $\sigma(\theta)$ is some $\theta_i$; and $\sigma$ is determined by $\sigma(\theta)$, since every element of $K$ is
$\sum a_k\theta^k$ with $a_k\in\Q$. Conversely, for each $i$ the map
$\Q[X]/(m_\theta)\to\C$, $P\mapsto P(\theta_i)$, is a well-defined ring homomorphism (as $m_\theta(\theta_i)=0$), and
composing it with $K\cong\Q[X]/(m_\theta)$ gives an embedding with $\theta\mapsto\theta_i$.
\end{proof}

An embedding is \emph{real} if $\sigma(K)\subseteq\R$, and \emph{complex} otherwise. If $\sigma$ is complex, so is its
conjugate $\bar\sigma:x\mapsto\overline{\sigma(x)}$, and $\bar\sigma\ne\sigma$; so complex embeddings come in pairs.
We write $r_1$ for the number of real embeddings and $r_2$ for the number of pairs, so that
\[
n=r_1+2r_2.
\]
The pair $(r_1,r_2)$ is the \emph{signature} of $K$. For $\Q(\sqrt2)$ it is $(2,0)$, for $\Q(i)$ it is $(0,1)$, and
for $\Q(\sqrt[3]2)$ it is $(1,1)$: the roots of $X^3-2$ are $\sqrt[3]2$, $\sqrt[3]2\,\omega$ and $\sqrt[3]2\,\bar\omega$
with $\omega=e^{2i\pi/3}$ (Figure~\ref{fig:embeddings}). A field with $r_2=0$ is \emph{totally real}, one with $r_1=0$
\emph{totally complex}. The cyclotomic field $\Q(\zeta_m)$, $m\ge3$, is totally complex, with $r_2=\varphi(m)/2$.

\begin{figure}[ht]
\centering
\fig{ch01-embeddings}
\caption{The three embeddings of $K=\Q(\sqrt[3]2)=\Q(\theta)$ send $\theta$ to the three cube roots of $2$. One is real,
the other two are exchanged by complex conjugation: the signature is $(r_1,r_2)=(1,1)$. The field $K$ itself is a
subfield of $\R$, but it ``sees'' the two complex roots through its embeddings.}
\label{fig:embeddings}
\end{figure}

\begin{think}
The field $\Q(\sqrt[3]2)$ is a subfield of $\R$. How can it have complex embeddings? And is $\Q(\sqrt[3]2)$ equal to
$\Q(\sqrt[3]2\,\omega)$, or only isomorphic to it?
\end{think}

The embeddings give each element $\alpha\in K$ its \emph{conjugates} $\sigma_1(\alpha),\dots,\sigma_n(\alpha)$. They
are the roots of a polynomial that we can compute without knowing the $\sigma_i$.

\begin{proposition}\label{prop:charpoly}
For $\alpha\in K$, let $\mu_\alpha:K\to K$, $x\mapsto\alpha x$, a $\Q$-linear map, and
$\chi_\alpha=\prod_{i=1}^n(X-\sigma_i(\alpha))$. Then
\[
\chi_\alpha=\det(X-\mu_\alpha)=m_\alpha^{\,n/d},\qquad d=[\Q(\alpha):\Q].
\]
In particular $\chi_\alpha\in\Q[X]$, and each conjugate of $\alpha$ is a root of $m_\alpha$, repeated $n/d$ times
among the $\sigma_i(\alpha)$.
\end{proposition}
\begin{proof}
Let $F=\Q(\alpha)$, $e=n/d=[K:F]$ (tower law). Take a basis $\beta_1,\dots,\beta_e$ of $K$ over $F$. Then the
$\alpha^k\beta_j$ ($0\le k<d$, $1\le j\le e$) form a basis of $K$ over $\Q$ (proof of the tower law). On the span of
$\beta_j,\alpha\beta_j,\dots,\alpha^{d-1}\beta_j$, multiplication by $\alpha$ acts by the companion matrix of
$m_\alpha$, whose characteristic polynomial is $m_\alpha$. So the matrix of $\mu_\alpha$ is block diagonal with $e$
copies of it, and $\det(X-\mu_\alpha)=m_\alpha^e$.

Now count embeddings. Each embedding $\sigma$ of $K$ restricts to an embedding of $F$, which sends $\alpha$ to a root of
$m_\alpha$; there are $d$ embeddings $\tau$ of $F$, one for each root (Proposition~\ref{prop:embeddings}). We claim each
$\tau$ is the restriction of exactly $e$ embeddings of $K$. Write $K=F(\gamma)$ (primitive element for $K/F$, same
exercise), with minimal polynomial $g$ over $F$ of degree $e$. An embedding of $K$ extending $\tau$ is determined by the
image of $\gamma$, which must be a root of $\tau g$ (the polynomial with $\tau$ applied to the coefficients); $\tau g$
is irreducible over $\tau(F)$ (a factorisation would pull back by $\tau^{-1}$), so it has $e$ distinct roots, and each
gives an extension (Tome~2, Lemma~8.12). So $\chi_\alpha=\prod_\tau(X-\tau(\alpha))^e=m_\alpha^e$, since the
$\tau(\alpha)$ are the $d$ roots of $m_\alpha$.
\end{proof}

\section{Algebraic integers}

\begin{definition}
A complex number $\alpha$ is an \emph{algebraic integer} if it is a root of a monic polynomial with coefficients in $\Z$.
For a number field $K$, $\mathcal O_K$ denotes the set of algebraic integers in $K$: the \emph{integers} of $K$.
\end{definition}

So $\sqrt2$, $\frac{1+\sqrt5}2$ (root of $X^2-X-1$), $e^{2i\pi/7}$ (root of $X^7-1$) and $\sqrt[3]2$ are algebraic
integers; $\frac12$ is not (rational root theorem). Is $\sqrt2+\sqrt[3]2$ one? Finding a monic equation is a tedious
elimination. The following criterion replaces equations by finitely generated groups, and makes the answer
immediate.

\begin{proposition}\label{prop:intcrit}
For $\alpha\in\C$ the following are equivalent:
\begin{enumerate}[label=(\alph*)]
\item $\alpha$ is an algebraic integer;
\item the additive group $\Z[\alpha]$ is finitely generated;
\item there is a non-zero finitely generated additive subgroup $M\subseteq\C$ with $\alpha M\subseteq M$.
\end{enumerate}
\end{proposition}
\begin{proof}
(a)$\Rightarrow$(b). If $\alpha^d=c_{d-1}\alpha^{d-1}+\dots+c_0$ with $c_k\in\Z$, then by induction every power
$\alpha^m$ is an integer combination of $1,\alpha,\dots,\alpha^{d-1}$, which therefore generate $\Z[\alpha]$.
(b)$\Rightarrow$(c). Take $M=\Z[\alpha]$, which contains $1\ne0$.
(c)$\Rightarrow$(a). Let $m_1,\dots,m_k$ generate $M$. Then $\alpha m_i=\sum_ja_{ij}m_j$ with $a_{ij}\in\Z$, that is,
$(\alpha I-A)v=0$ for $A=(a_{ij})$ and the vector $v=(m_1,\dots,m_k)^{\mathsf T}\in\C^k$, which is non-zero since
$M\ne0$. So $\det(\alpha I-A)=0$, and $\det(XI-A)$ is a monic polynomial with integer coefficients vanishing at $\alpha$.
\end{proof}

\begin{corollary}\label{cor:ring}
Sums and products of algebraic integers are algebraic integers. In particular $\mathcal O_K$ is a subring of $K$.
\end{corollary}
\begin{proof}
If $\alpha^d$ and $\beta^e$ are integer combinations of lower powers, the group $M=\Z[\alpha,\beta]$ is generated by
the $de$ products $\alpha^i\beta^j$ with $i<d$, $j<e$. It is stable under multiplication by $\alpha+\beta$ and by
$\alpha\beta$, so both are algebraic integers by (c).
\end{proof}

So $\sqrt2+\sqrt[3]2$ is an algebraic integer, without computing its minimal polynomial (which is
$X^6-6X^4-4X^3+12X^2-24X-4$). The same argument, with more generators, proves transitivity.

\begin{corollary}\label{cor:transitive}
If $\alpha\in\C$ is a root of a monic polynomial whose coefficients $b_0,\dots,b_{d-1}$ are algebraic integers, then
$\alpha$ is an algebraic integer.
\end{corollary}
\begin{proof}
Each $b_k$ satisfies a monic integer equation of some degree $e_k$. The group $M$ generated by the monomials
$b_0^{i_0}\cdots b_{d-1}^{i_{d-1}}\alpha^j$ with $i_k<e_k$ and $j<d$ is finitely generated and stable under
multiplication by $\alpha$: in $\alpha\cdot b^{i}\alpha^{d-1}$, replace $\alpha^d$ by $-\sum b_k\alpha^k$, then reduce
each power $b_k^{e_k}$ in the same way. By Proposition~\ref{prop:intcrit}(c), $\alpha$ is an algebraic integer.
\end{proof}

\begin{proposition}\label{prop:minpolyZ}
Let $\alpha$ be an algebraic number. Then $\alpha$ is an algebraic integer if and only if $m_\alpha\in\Z[X]$. In
particular $\mathcal O_K\cap\Q=\Z$, the conjugates of an algebraic integer are algebraic integers, and for
$\alpha\in\mathcal O_K$ the polynomial $\chi_\alpha$ has integer coefficients.
\end{proposition}
\begin{proof}
If $m_\alpha\in\Z[X]$, $\alpha$ is an algebraic integer by definition. Conversely let $P\in\Z[X]$ be monic with
$P(\alpha)=0$. Then $m_\alpha\mid P$ in $\Q[X]$, so every root $\alpha'$ of $m_\alpha$ is a root of $P$, hence an
algebraic integer. The coefficients of $m_\alpha=\prod(X-\alpha')$ are, up to sign, sums of products of the $\alpha'$,
so they are algebraic integers (Corollary~\ref{cor:ring}) and rational: they are integers, because a rational
algebraic integer is a root of a monic integer polynomial and hence in $\Z$ (rational root theorem). This proves the
equivalence, and the other statements follow: a rational $a$ has $m_a=X-a$; the conjugates of $\alpha$ are roots of the
same $m_\alpha$; and $\chi_\alpha=m_\alpha^{n/d}$ (Proposition~\ref{prop:charpoly}).
\end{proof}

\begin{proposition}\label{prop:denominators}
Every $\alpha\in K$ has a multiple $c\alpha\in\mathcal O_K$ with $c\in\Z$, $c\ge1$. Hence $K$ is the field of
fractions of $\mathcal O_K$, and $\mathcal O_K$ contains a basis of $K$ over $\Q$.
\end{proposition}
\begin{proof}
Let $a_d\alpha^d+\dots+a_0=0$ with $a_k\in\Z$, $a_d\ne0$ (clear the denominators of $m_\alpha$). Multiplying by
$a_d^{d-1}$, $(a_d\alpha)^d+a_{d-1}(a_d\alpha)^{d-1}+a_{d-2}a_d(a_d\alpha)^{d-2}+\dots+a_0a_d^{d-1}=0$, a monic integer
equation for $a_d\alpha$. Apply this to a basis of $K$ over $\Q$.
\end{proof}

\section{Norm, trace and discriminant}

\begin{definition}
The \emph{trace} and the \emph{norm} of $\alpha\in K$ are
\[
\operatorname{Tr}_{K/\Q}(\alpha)=\sum_{i=1}^n\sigma_i(\alpha)=\operatorname{tr}\mu_\alpha,\qquad
N_{K/\Q}(\alpha)=\prod_{i=1}^n\sigma_i(\alpha)=\det\mu_\alpha,
\]
the equalities coming from Proposition~\ref{prop:charpoly} (the trace and the determinant of a matrix are, up to sign,
the coefficients of $X^{n-1}$ and $X^0$ in its characteristic polynomial).
\end{definition}

We drop the subscript $K/\Q$ when the field is clear. The second description shows that $\operatorname{Tr}$ and $N$
take rational values, that $\operatorname{Tr}$ is $\Q$-linear and $N$ is multiplicative ($\mu_{\alpha\beta}=
\mu_\alpha\mu_\beta$), and that $N(a)=a^n$, $\operatorname{Tr}(a)=na$ for $a\in\Q$. By
Proposition~\ref{prop:minpolyZ}, they take integer values on $\mathcal O_K$. For $K=\Q(\sqrt d)$ we recover Tome~1's
$\operatorname{Tr}(a+b\sqrt d)=2a$ and $N(a+b\sqrt d)=a^2-db^2$. For $K=\Q(\sqrt[3]2)$ and $\alpha=a+b\theta+c\theta^2$,
the matrix of $\mu_\alpha$ in the basis $1,\theta,\theta^2$ gives
\[
N(a+b\sqrt[3]2+c\sqrt[3]4)=\det\begin{pmatrix}a&2c&2b\\b&a&2c\\c&b&a\end{pmatrix}=a^3+2b^3+4c^3-6abc.
\]

\begin{proposition}\label{prop:units}
An element $\alpha\in\mathcal O_K$ is a unit of $\mathcal O_K$ if and only if $N(\alpha)=\pm1$.
\end{proposition}
\begin{proof}
If $\alpha\beta=1$ with $\beta\in\mathcal O_K$, then $N(\alpha)N(\beta)=1$ with two integers, so $N(\alpha)=\pm1$.
Conversely, suppose $N(\alpha)=\pm1$, and number the embeddings so that $\sigma_1$ is the inclusion $K\subseteq\C$.
Then $\alpha\ne0$, and $\gamma=\prod_{i\ge2}\sigma_i(\alpha)$ is a product of algebraic integers, hence an algebraic
integer. It lies in $K$, because $\gamma=N(\alpha)/\alpha$. So $\gamma\in\mathcal O_K$, and $\alpha\cdot N(\alpha)\gamma=
N(\alpha)^2=1$: the inverse of $\alpha$ is $N(\alpha)\gamma\in\mathcal O_K$.
\end{proof}

For example $N(\sqrt[3]2-1)=-1+2=1$ by the formula above, so $\sqrt[3]2-1$ is a unit of $\Z[\sqrt[3]2]$; its inverse
is $1+\sqrt[3]2+\sqrt[3]4$, as $(\theta-1)(\theta^2+\theta+1)=\theta^3-1=1$. Its powers give infinitely many integer
solutions of the cubic equation $a^3+2b^3+4c^3-6abc=1$, the cubic analogue of Pell's equation. Chapter~4 shows that
they are all the units, up to sign.

The basic invariant of a family of $n$ elements of $K$ measures how far they are from being a basis.

\begin{definition}
The \emph{discriminant} of $\omega_1,\dots,\omega_n\in K$ is
\[
\operatorname{disc}(\omega_1,\dots,\omega_n)=\det\bigl(\sigma_i(\omega_j)\bigr)_{1\le i,j\le n}^2.
\]
\end{definition}

The square makes the definition independent of the order of the embeddings and of the $\omega_j$.

\begin{proposition}\label{prop:disc}
Let $\omega_1,\dots,\omega_n\in K$.
\begin{enumerate}[label=(\alph*)]
\item $\operatorname{disc}(\omega)=\det\bigl(\operatorname{Tr}(\omega_i\omega_j)\bigr)$; so it is rational, and an
integer if the $\omega_j$ are in $\mathcal O_K$.
\item If $\omega'_i=\sum_jc_{ij}\omega_j$ with $C=(c_{ij})\in M_n(\Q)$, then
$\operatorname{disc}(\omega')=(\det C)^2\operatorname{disc}(\omega)$.
\item If $K=\Q(\theta)$, with $\theta_i=\sigma_i(\theta)$ and $f=m_\theta$, then
\[
\operatorname{disc}(1,\theta,\dots,\theta^{n-1})=\prod_{i<j}(\theta_i-\theta_j)^2
=(-1)^{n(n-1)/2}N\bigl(f'(\theta)\bigr)\ne0.
\]
\item $\operatorname{disc}(\omega)\ne0$ if and only if $\omega_1,\dots,\omega_n$ is a basis of $K$ over $\Q$.
\end{enumerate}
\end{proposition}
\begin{proof}
(a) With $A=(\sigma_i(\omega_j))$, the $(j,k)$ entry of $A^{\mathsf T}A$ is
$\sum_i\sigma_i(\omega_j)\sigma_i(\omega_k)=\sum_i\sigma_i(\omega_j\omega_k)=\operatorname{Tr}(\omega_j\omega_k)$, and
$\det(A^{\mathsf T}A)=(\det A)^2$. (b) The matrix $(\sigma_i(\omega'_j))$ is $(\sigma_i(\omega_k))C^{\mathsf T}$,
because $\sigma_i$ fixes the rational $c_{jk}$. (c) The matrix $(\theta_i^{\,j-1})$ is a Vandermonde matrix, with
determinant $\prod_{i<j}(\theta_j-\theta_i)$. Then $f=\prod_j(X-\theta_j)$ gives $f'(\theta_i)=\prod_{j\ne i}(\theta_i-
\theta_j)$, and $N(f'(\theta))=\prod_i f'(\theta_i)=\prod_{i\ne j}(\theta_i-\theta_j)$, which is the product over pairs
$i<j$ of $(\theta_i-\theta_j)(\theta_j-\theta_i)=-(\theta_i-\theta_j)^2$; there are $n(n-1)/2$ pairs. The roots are
distinct, so this is non-zero. (d) A basis is obtained from $1,\theta,\dots,\theta^{n-1}$ by an invertible $C$, so its
discriminant is non-zero by (b) and (c). If the $\omega_j$ are dependent, some non-zero rational combination of the
columns of $(\sigma_i(\omega_j))$ vanishes, and the determinant is $0$.
\end{proof}

The number $\prod_{i<j}(\theta_i-\theta_j)^2$ is the \emph{discriminant} $\operatorname{disc}(f)$ of the polynomial
$f$. For $X^2+bX+c$ it is $b^2-4c$; for $X^3+aX+b$ it is $-4a^3-27b^2$ (Exercise~\ref{ex:disccubic}), so
$\operatorname{disc}(X^3-2)=-108$, $\operatorname{disc}(X^3-X-1)=-23$ and $\operatorname{disc}(X^3-3X+1)=81$.

\section{Integral bases and the discriminant of a field}

\begin{theorem}\label{thm:intbasis}
Let $K$ be a number field of degree $n$. Then $\mathcal O_K$ is a free abelian group of rank $n$: there are
$\omega_1,\dots,\omega_n\in\mathcal O_K$ such that every $\alpha\in\mathcal O_K$ is uniquely $\sum a_j\omega_j$ with
$a_j\in\Z$. Such an $(\omega_j)$ is an \emph{integral basis}; it is also a basis of $K$ over $\Q$.
\end{theorem}
\begin{proof}
By Proposition~\ref{prop:denominators}, $\mathcal O_K$ contains a $\Q$-basis $\beta_1,\dots,\beta_n$ of $K$. Let
$D=\operatorname{disc}(\beta)$, a non-zero integer (Proposition~\ref{prop:disc}). For $\alpha\in\mathcal O_K$ write
$\alpha=\sum_jc_j\beta_j$ with $c_j\in\Q$. Multiplying by $\beta_i$ and taking traces,
\[
\operatorname{Tr}(\alpha\beta_i)=\sum_jc_j\operatorname{Tr}(\beta_i\beta_j),\qquad i=1,\dots,n,
\]
a linear system with integer right-hand sides and an integer matrix $T=(\operatorname{Tr}(\beta_i\beta_j))$ of
determinant $D$. By Cramer's rule, $c_j=D_j/D$ with $D_j\in\Z$. Hence
\[
\Z\beta_1+\dots+\Z\beta_n\ \subseteq\ \mathcal O_K\ \subseteq\ \tfrac1D\bigl(\Z\beta_1+\dots+\Z\beta_n\bigr).
\]
The right-hand group is free of rank $n$, with basis $\beta_j/D$. As a subgroup of it, $\mathcal O_K$ is free of rank
at most $n$ (Tome~2, Lemma~7.5), and its rank is at least $n$ because it contains $n$ independent elements. A $\Z$-basis
of $\mathcal O_K$ spans $K$ over $\Q$ (every element of $K$ is a rational multiple of an element of $\mathcal O_K$), so
it is a $\Q$-basis.
\end{proof}

If $(\omega_j)$ and $(\omega'_j)$ are two integral bases, each is an integer combination of the other, so the change of
basis matrix $C$ and its inverse have integer entries, $\det C=\pm1$, and the two discriminants are equal.

\begin{definition}
The \emph{discriminant} $d_K$ of $K$ is the discriminant of any integral basis. It is a non-zero integer.
\end{definition}

For a quadratic field this is $d_K=\operatorname{disc}(1,\sqrt d)=4d$ if $d\equiv2,3\pmod4$ and
$\operatorname{disc}(1,\frac{1+\sqrt d}2)=d$ if $d\equiv1\pmod4$, by Tome~1's integral bases. Figure~\ref{fig:mink}
shows the meaning of $\sqrt{|d_K|}$: send $\alpha\in\mathcal O_K$ to the point $(\sigma_1(\alpha),\dots,\sigma_n(\alpha))$
of $\R^n$ (for a totally real field); then $\mathcal O_K$ becomes a lattice whose fundamental parallelotope has volume
$|\det(\sigma_i(\omega_j))|=\sqrt{|d_K|}$. The lattice picture is the subject of Chapter~3.

\begin{figure}[ht]
\centering
\fig{ch01-minkowski}
\caption{$\Z[\sqrt2]$ seen through its two embeddings: $a+b\sqrt2\mapsto(a+b\sqrt2,\,a-b\sqrt2)$. The ring becomes a
lattice in the plane, with fundamental parallelogram of area $|\det\begin{psmallmatrix}1&\sqrt2\\1&-\sqrt2\end{psmallmatrix}|=
2\sqrt2=\sqrt{d_K}$, $d_K=8$. The units are the lattice points on the hyperbolas $xy=\pm1$ (larger dots); they march
off to infinity along the hyperbolas, by powers of $1+\sqrt2$.}
\label{fig:mink}
\end{figure}

\begin{proposition}[the sign of the discriminant]\label{prop:sign}
The sign of $d_K$ is $(-1)^{r_2}$.
\end{proposition}
\begin{proof}
Let $\Delta=\det(\sigma_i(\omega_j))$ for an integral basis. Complex conjugation applied to the matrix fixes the rows of
the real embeddings and exchanges the two rows of each complex pair; so $\bar\Delta=(-1)^{r_2}\Delta$. If $r_2$ is even,
$\Delta$ is real and $d_K=\Delta^2>0$; if $r_2$ is odd, $\Delta$ is purely imaginary and $d_K=\Delta^2<0$.
\end{proof}

\begin{proposition}[the index formula]\label{prop:index}
Let $\alpha_1,\dots,\alpha_n\in\mathcal O_K$ be linearly independent over $\Q$, and $M=\Z\alpha_1+\dots+\Z\alpha_n$.
Then $M$ has finite index in $\mathcal O_K$ and
\[
\operatorname{disc}(\alpha_1,\dots,\alpha_n)=[\mathcal O_K:M]^2\,d_K.
\]
In particular, if $K=\Q(\theta)$ with $\theta\in\mathcal O_K$, then $\operatorname{disc}(m_\theta)=
[\mathcal O_K:\Z[\theta]]^2\,d_K$; and if $\operatorname{disc}(m_\theta)$ is squarefree, then
$\mathcal O_K=\Z[\theta]$.
\end{proposition}
\begin{proof}
Write $\alpha_i=\sum_jc_{ij}\omega_j$ in an integral basis, with $C=(c_{ij})\in M_n(\Z)$, invertible over $\Q$ since the
$\alpha_i$ are independent. Through the basis $(\omega_j)$, $\mathcal O_K\cong\Z^n$ and $M$ corresponds to
$C^{\mathsf T}\Z^n$, of index $|\det C|$ (Tome~2, Corollary~7.8). Proposition~\ref{prop:disc}(b) gives
$\operatorname{disc}(\alpha)=(\det C)^2d_K$. For $M=\Z[\theta]$, with basis $1,\theta,\dots,\theta^{n-1}$, the discriminant
is $\operatorname{disc}(m_\theta)$ by Proposition~\ref{prop:disc}(c). If it is squarefree, the square
$[\mathcal O_K:\Z[\theta]]^2$ dividing it is $1$.
\end{proof}

The number $[\mathcal O_K:\Z[\theta]]$ is the \emph{index} of $\theta$. For example $X^3-X-1$ has discriminant $-23$,
squarefree, so if $\theta^3=\theta+1$ then $\mathcal O_K=\Z[\theta]$ and $d_K=-23$. Two different defining polynomials
of the same field can have different discriminants, but they all are $d_K$ times a square.

\begin{think}
The discriminants of the quadratic fields are $5, 8, 12, 13, 17, 21, 24, \dots$ and $-3,-4,-7,-8,-11,-15,\dots$. Do you
see a pattern modulo $4$? Guess what remainders modulo $4$ the discriminant of a cubic field can have.
\end{think}

\section{Stickelberger's theorem}

\begin{theorem}[Stickelberger, 1897]\label{thm:stickelberger}
For every number field $K$, $d_K\equiv0$ or $1\pmod 4$.
\end{theorem}
\begin{proof}
Let $(\omega_j)$ be an integral basis and $\Delta=\det(\sigma_i(\omega_j))=\sum_\pi\operatorname{sgn}(\pi)\prod_i
\sigma_i(\omega_{\pi(i)})$, the sum over permutations $\pi$ of $\{1,\dots,n\}$. Split it as $\Delta=P-N$, where $P$
collects the terms with $\pi$ even and $N$ those with $\pi$ odd (without their sign). Both are algebraic integers.

Let $L\subseteq\C$ be the splitting field of $m_\theta$ for a primitive element $\theta$ of $K$: it contains every
$\sigma_i(K)$, and it is Galois over $\Q$. For an automorphism $\tau$ of $L$, the maps $\tau\circ\sigma_i$ are again the
$n$ embeddings of $K$, in some order: $\tau\circ\sigma_i=\sigma_{\rho(i)}$ for a permutation $\rho$. Applying $\tau$ to
$\Delta$ permutes the rows of the matrix by $\rho$. If $\rho$ is even, $\tau(P)=P$ and $\tau(N)=N$; if $\rho$ is odd,
$\tau$ exchanges $P$ and $N$. In both cases $\tau$ fixes $P+N$ and $PN$, which are therefore rational (Galois
theory), and algebraic integers: integers. Finally
\[
d_K=(P-N)^2=(P+N)^2-4PN\equiv(P+N)^2\pmod 4,
\]
and a square is $0$ or $1$ modulo $4$.
\end{proof}

So no number field has discriminant $\equiv2,3\pmod4$. For instance $X^3+X+1$ has discriminant $-31\equiv1\pmod4$,
squarefree, so it defines a cubic field with $d_K=-31$; and no cubic field has discriminant $-30$, $-26$ or $-22$,
whatever polynomial one tries. For quadratic fields the theorem is visible in Tome~1's list: $d_K$ is $d\equiv1$ or
$4d\equiv0\pmod 4$. Stickelberger's theorem is one of the first constraints on which integers can be discriminants;
Chapter~3 adds a lower bound on $|d_K|$ in terms of $n$.

\section{Computing rings of integers}

The index formula bounds the index: $[\mathcal O_K:\Z[\theta]]^2$ divides $\operatorname{disc}(m_\theta)$. So only primes
$p$ with $p^2\mid\operatorname{disc}(m_\theta)$ can divide the index, and for each one we need a test. Eisenstein's
criterion provides one.

\begin{proposition}\label{prop:eisenstein}
Let $\theta\in\mathcal O_K$, $K=\Q(\theta)$, and suppose $m_\theta=X^n+a_{n-1}X^{n-1}+\dots+a_0$ is
\emph{$p$-Eisenstein}: $p\mid a_i$ for all $i$, and $p^2\nmid a_0$. Then $p$ does not divide
$[\mathcal O_K:\Z[\theta]]$.
\end{proposition}
\begin{proof}
Suppose $p$ divides the index. The finite group $\mathcal O_K/\Z[\theta]$ then has an element of order $p$ (Cauchy's
theorem): there is $\beta\in\mathcal O_K\setminus\Z[\theta]$ with $p\beta\in\Z[\theta]$. Write $p\beta=\sum_{i<n}b_i
\theta^i$ with $b_i\in\Z$, not all divisible by $p$ (else $\beta\in\Z[\theta]$), and let $r$ be the smallest index with
$p\nmid b_r$. Subtracting from $\beta$ the element $\sum_{i<r}\frac{b_i}p\theta^i$ of $\Z[\theta]$, we get
$\gamma\in\mathcal O_K$ with
\[
p\gamma=b_r\theta^r+b_{r+1}\theta^{r+1}+\dots+b_{n-1}\theta^{n-1}.
\]
Since $m_\theta$ is Eisenstein, $\theta^n=-(a_{n-1}\theta^{n-1}+\dots+a_0)\in p\Z[\theta]$, so $\theta^n/p\in
\Z[\theta]$, and $\theta^k/p\in\Z[\theta]$ for every $k\ge n$. Multiply the display by $\theta^{n-1-r}/p$:
\[
\frac{b_r\theta^{n-1}}p=\gamma\theta^{n-1-r}-\sum_{i>r}b_i\frac{\theta^{\,n-1-r+i}}p,
\]
and every term on the right is in $\mathcal O_K$ (the exponents $n-1-r+i$ are $\ge n$). So
$b_r\theta^{n-1}/p\in\mathcal O_K$ and its norm is an integer. But $N(\theta)=(-1)^na_0$, so
\[
N\!\left(\frac{b_r\theta^{n-1}}p\right)=\frac{b_r^n\,(\pm a_0)^{n-1}}{p^n},
\]
whose numerator is divisible by exactly $p^{n-1}$ (as $p\nmid b_r$ and $p\,\|\,a_0$): not an integer, a contradiction.
\end{proof}

Since $\Z[\theta]=\Z[\theta+c]$ for every $c\in\Z$, one may apply the criterion to a translate of $\theta$. Here are the
two fields of the opening question.

\begin{example}[{$\Q(\sqrt[3]2)$}]\label{exa:cuberoot2}
Let $\theta=\sqrt[3]2$. Then $\operatorname{disc}(X^3-2)=-4\cdot0-27\cdot4=-108=-2^2\cdot3^3$, so the index divides $6$.
$X^3-2$ is $2$-Eisenstein, so $2$ does not divide the index. For $3$, $\theta+1$ has minimal polynomial
$(X-1)^3-2=X^3-3X^2+3X-3$, which is $3$-Eisenstein, so $3$ does not divide the index of $\theta+1$, which is that of
$\theta$. Hence $\mathcal O_K=\Z[\sqrt[3]2]$ and $d_K=-108$.
\end{example}

\begin{example}[{$\Q(\sqrt[3]{10})$}]\label{exa:cuberoot10}
Let $\theta=\sqrt[3]{10}$ and $\beta=\frac{1+\theta+\theta^2}3$. Then $\beta$ is an algebraic integer. Indeed
$(1+\theta+\theta^2)(\theta-1)=\theta^3-1=9$, so $y=\theta-1=\frac3\beta$, and $y$ is a root of
$(X+1)^3-10=X^3+3X^2+3X-9$. Substituting $y=3/\beta$ and multiplying by $\beta^3/9$ gives
\[
\beta^3-\beta^2-3\beta-3=0,
\]
a monic equation with integer coefficients. So $\beta\in\mathcal O_K\setminus\Z[\theta]$, and $3$ divides the index. Exercise~\ref{ex:cuberoot10} shows that $\mathcal O_K=\Z[\theta,\beta]$, of index $3$, and
$d_K=-2700/9=-300$.
\end{example}

The difference between $2$ and $10$ is a congruence: for a squarefree integer $m\ne0,\pm1$, $\mathcal O_{\Q(\sqrt[3]m)}=\Z[\sqrt[3]m]$
when $m\not\equiv\pm1\pmod9$, and the index is $3$ when $m\equiv\pm1\pmod 9$, because then
$X^3-m\equiv(X\mp1)^3\pmod 9$ has a root of unusually high multiplicity modulo $3$. Chapter~2 explains why such
repeated roots modulo $p$ are exactly what makes $p$ divide the index.

\begin{think}
Is there a number field $K$ in which $\mathcal O_K\ne\Z[\theta]$ for \emph{every} $\theta\in\mathcal O_K$? Dedekind found one
in 1878 (Exercise~\ref{ex:dedekind}); the reason, which needs Chapter~2, is that a prime can split into more prime
ideals of degree $1$ than there are elements in the residue field.
\end{think}

\section{Exercises}

\begin{exercise}[$\star$]\label{ex:quadratic}
Recover Tome~1's description of the integers of $\Q(\sqrt d)$, $d$ squarefree, from this chapter: compute
$\operatorname{disc}(1,\sqrt d)$, use the index formula to see that $[\mathcal O_K:\Z[\sqrt d]]\in\{1,2\}$, and decide
when $\frac{a+b\sqrt d}2$ ($a,b$ odd) is integral by computing its minimal polynomial.
\end{exercise}

\begin{exercise}[$\star$]\label{ex:roots}
Let $\alpha\in\C$ and $k\ge1$ with $\alpha^k$ an algebraic integer. Show that $\alpha$ is an algebraic integer. Is
$\sqrt{1+\sqrt2}+\sqrt[5]{3}$ an algebraic integer?
\end{exercise}

\begin{exercise}[$\star$]\label{ex:disccubic}
Show that $\operatorname{disc}(X^3+aX+b)=-4a^3-27b^2$. Deduce that if $\theta^3=\theta+1$, then
$\mathcal O_{\Q(\theta)}=\Z[\theta]$ with discriminant $-23$, and that $\Q(\theta)$ has exactly one real embedding.
\end{exercise}

\begin{exercise}[$\star\star$]\label{ex:normform}
\emph{The cubic Pell equation.} Show that $N(a+b\sqrt[3]2+c\sqrt[3]4)=a^3+2b^3+4c^3-6abc$, and that the equation
$a^3+2b^3+4c^3-6abc=1$ has infinitely many solutions in integers. Find the one coming from $(\sqrt[3]2-1)^3$.
\end{exercise}

\begin{exercise}[$\star\star$]\label{ex:discbinom}
Show that $\operatorname{disc}(X^n-a)=(-1)^{(n-1)(n-2)/2}\,n^n\,a^{n-1}$, and more generally that
$\operatorname{disc}(f)=(-1)^{n(n-1)/2}\prod_if'(\theta_i)$ for monic $f$ with roots $\theta_i$. Check the formula on
$X^3-2$.
\end{exercise}

\begin{exercise}[$\star\star$]\label{ex:kronecker}
\emph{Kronecker's theorem.} Let $\alpha$ be an algebraic integer all of whose conjugates have absolute value at most $1$.
Show that $\alpha=0$ or $\alpha$ is a root of unity. \emph{(Look at the characteristic polynomials of the powers of
$\alpha$.)} Is the conclusion still true for an algebraic number that is not an integer?
\end{exercise}

\begin{exercise}[$\star\star$]\label{ex:sqrt23}
Let $\alpha=\sqrt2+\sqrt3$. Show that $m_\alpha=X^4-10X^2+1$ and that $\Q(\alpha)=\Q(\sqrt2,\sqrt3)$. Show that
$\frac{\sqrt2+\sqrt6}2$ is an algebraic integer not in $\Z[\sqrt2,\sqrt3]$.
\end{exercise}

\begin{exercise}[$\star\star$]\label{ex:biquadratic}
\emph{The integers of $\Q(\sqrt2,\sqrt3)$.} Show that $1,\sqrt2,\sqrt3,\frac{\sqrt2+\sqrt6}2$ is an integral basis and
that $d_K=2304=2^8\cdot3^2$. Deduce that $[\mathcal O_K:\Z[\sqrt2+\sqrt3]]=8$. \emph{(Write $x=a+b\sqrt2+c\sqrt3+
d\sqrt6$; the traces of $x$, $x\sqrt2$, $x\sqrt3$ and $x\sqrt6$ to the three quadratic subfields must be integers.)}
\end{exercise}

\begin{exercise}[$\star\star$]\label{ex:cuberoot10}
Complete Example~\ref{exa:cuberoot10}: show that $1,\theta,\beta$ is an integral basis of $\Q(\sqrt[3]{10})$ and that
$d_K=-300$. \emph{(Use the index formula and Eisenstein at $2$ and $5$.)}
\end{exercise}

\begin{exercise}[$\star\star$]\label{ex:totallyreal}
Let $\theta=2\cos\frac{2\pi}9$. Show that $\theta^3-3\theta+1=0$, that $\Q(\theta)$ is totally real with
$\mathcal O_K=\Z[\theta]$ and $d_K=81$, and that the three roots of $X^3-3X+1$ are $\theta$, $\theta^2-2$ and
$2-\theta-\theta^2$, so that $\Q(\theta)$ is Galois over $\Q$.
\end{exercise}

\begin{exercise}[$\star\star\star$]\label{ex:dedekind}
\emph{Dedekind's cubic field.} Let $\theta$ be a root of $f=X^3-X^2-2X-8$ (irreducible: it has no integer root). Show
that $\operatorname{disc}(f)=-2012=-4\cdot503$, that $\beta=\frac{\theta^2+\theta}2$ is an algebraic integer, and that
$1,\theta,\beta$ is an integral basis, so that $d_K=-503$.
\end{exercise}

\begin{exercise}[$\star\star\star$]\label{ex:traceform}
\emph{The signature of the trace form.} Show that the quadratic form $Q(x)=\operatorname{Tr}(x^2)$ on the real vector
space $K\otimes\R\cong\R^n$ (that is, with real coordinates in a $\Q$-basis) has signature $(r_1+r_2,r_2)$. Deduce
another proof that the sign of $d_K$ is $(-1)^{r_2}$, and that $K$ is totally real if and only if $\operatorname{Tr}(x^2)>0$
for all $x\in K^\times$.
\end{exercise}

\begin{beauty}
The whole chapter rests on one idea, Proposition~\ref{prop:intcrit}: an element is integral when it acts on a lattice.
It turns equations into linear algebra, and makes it obvious that integers form a ring, which no one could check by
multiplying minimal polynomials. The discriminant is the volume of the lattice $\mathcal O_K$ (Figure~\ref{fig:mink}),
and Stickelberger's congruence is the first sign that this volume is not arbitrary. Both ideas reach far. The
discriminant of a field controls its ramified primes (Chapter~2), bounds its class number (Chapter~3), and appears in
its zeta function (Chapter~8); a theorem of Hermite says there are only finitely many fields of each discriminant.
Counting fields by discriminant is a modern subject in itself: Davenport and Heilbronn showed in 1971 that there are
about $\frac{x}{3\zeta(3)}$ cubic fields with $|d_K|<x$, and Bhargava extended this to quartic and quintic fields
(this tome's papers track, and Tome~14). Kronecker's theorem of Exercise~\ref{ex:kronecker} is the start of another story:
Lehmer asked in 1933 whether the largest conjugate of an algebraic integer that is not a root of unity can be
arbitrarily close to~$1$, in a precise sense; the question is still open.
\end{beauty}

\section{Hints}

Try each exercise for a while before reading its hint, and read the hint before the solution.

\begin{hint}{ex:quadratic}
$\operatorname{disc}(1,\sqrt d)=4d$, and $d$ is squarefree, so the index divides $2$. The element
$\frac{a+b\sqrt d}2$ has trace $a$ and norm $\frac{a^2-db^2}4$.
\end{hint}
\begin{hint}{ex:roots}
$\alpha$ is a root of $X^k-\alpha^k$; use Corollary~\ref{cor:transitive}.
\end{hint}
\begin{hint}{ex:disccubic}
Use $\operatorname{disc}(f)=-N(f'(\theta))$ for a cubic, with $f'(\theta)=3\theta^2+a$, and compute its norm as a
determinant, or with symmetric functions of the roots. The sign of the discriminant tells the signature.
\end{hint}
\begin{hint}{ex:normform}
Compute the matrix of multiplication by $a+b\theta+c\theta^2$ in the basis $1,\theta,\theta^2$. The powers of a unit of
norm $1$ are units of norm $1$.
\end{hint}
\begin{hint}{ex:discbinom}
$f'(\theta_i)=n\theta_i^{n-1}$, and $\prod_i\theta_i=(-1)^nf(0)$.
\end{hint}
\begin{hint}{ex:kronecker}
The characteristic polynomials of $\alpha^k$, $k\ge1$, have integer coefficients bounded in terms of $n$ only.
\end{hint}
\begin{hint}{ex:sqrt23}
Square $\alpha$ twice. For the second part, compute the square of $\frac{\sqrt2+\sqrt6}2$.
\end{hint}
\begin{hint}{ex:biquadratic}
Taking traces to $\Q(\sqrt2)$, $\Q(\sqrt3)$ and $\Q(\sqrt6)$ of an integer gives integers of these fields, whose
integral bases are known. Conclude with the parity of $a,b,c,d$ after multiplying by $2$.
\end{hint}
\begin{hint}{ex:cuberoot10}
$\operatorname{disc}(X^3-10)=-2^2\cdot3^3\cdot5^2$, so the index of $\Z[\theta]$ divides $30$. Then compare the chain
$\Z[\theta]\subseteq\Z+\Z\theta+\Z\beta\subseteq\mathcal O_K$ with the indices.
\end{hint}
\begin{hint}{ex:totallyreal}
Use $\cos3t=4\cos^3t-3\cos t$; for the ring of integers, translate by $1$ and use Eisenstein.
\end{hint}
\begin{hint}{ex:dedekind}
$\beta^2$ and $\theta\beta$ can be written in the basis $1,\theta,\beta$ with integer coefficients, so
$\Z[\theta,\beta]$ is stable under multiplication by $\beta$. For the index, $503$ is prime.
\end{hint}
\begin{hint}{ex:traceform}
Use the coordinates $x\mapsto(\sigma_i(x))$: $\operatorname{Tr}(x^2)=\sum_{\text{real}}\sigma(x)^2+
\sum_{\text{pairs}}2\operatorname{Re}(\sigma(x)^2)$, and $2\operatorname{Re}(z^2)=2(u^2-v^2)$ for $z=u+iv$.
\end{hint}

\section{Solutions}

\begin{solution}{ex:quadratic}
\idea The index formula reduces everything to the factor $4$ in $\operatorname{disc}(1,\sqrt d)=4d$: the only possible
extra integers have denominator $2$. Then decide which of them are integral from the trace and norm, as in Tome~1.
\writeup The embeddings send $\sqrt d$ to $\pm\sqrt d$, so $\operatorname{disc}(1,\sqrt d)=\det
\begin{psmallmatrix}1&\sqrt d\\1&-\sqrt d\end{psmallmatrix}^2=4d$. By the index formula,
$4d=[\mathcal O_K:\Z[\sqrt d]]^2d_K$, and since $d$ is squarefree the index is $1$ or $2$. If it is $2$, then
$2\mathcal O_K\subseteq\Z[\sqrt d]$, so every integer is $\frac{a+b\sqrt d}2$ with $a,b\in\Z$, and some such element
with $a$ or $b$ odd is integral. Its minimal polynomial is $X^2-aX+\frac{a^2-db^2}4$, which has integer coefficients
iff $a^2\equiv db^2\pmod4$. If $b$ is even, this forces $a$ even. If $b$ is odd, $db^2\equiv d\pmod 4$ and we need
$a^2\equiv d\pmod 4$, which is possible (with $a$ odd) exactly when $d\equiv1\pmod4$, and impossible for
$d\equiv2,3\pmod4$ (squares are $0,1$). So $\mathcal O_K=\Z[\sqrt d]$ with $d_K=4d$ when $d\equiv2,3\pmod4$, and
$\mathcal O_K=\Z[\frac{1+\sqrt d}2]$ with $d_K=d$ when $d\equiv1\pmod4$.
\end{solution}

\begin{solution}{ex:roots}
\idea A root of a monic polynomial with algebraic integer coefficients is an algebraic integer.
\writeup $\alpha$ is a root of $X^k-\alpha^k$, monic with coefficients $1,0,\dots,0,-\alpha^k$, all algebraic integers,
so $\alpha$ is an algebraic integer by Corollary~\ref{cor:transitive}. For the example: $\sqrt2$ is integral, so is
$1+\sqrt2$, then $\sqrt{1+\sqrt2}$ (its square is integral), and $\sqrt[5]3$; their sum is integral by
Corollary~\ref{cor:ring}.
\end{solution}

\begin{solution}{ex:disccubic}
\idea Use the formula
\[
\operatorname{disc}(f)=(-1)^{n(n-1)/2}\,N(f'(\theta)),
\]
which for $n=3$ is $-N(3\theta^2+a)$, and compute the
norm with symmetric functions of the roots $\theta_1,\theta_2,\theta_3$, which satisfy $\sum\theta_i=0$,
$\sum_{i<j}\theta_i\theta_j=a$, $\theta_1\theta_2\theta_3=-b$.
\writeup Assume first $a\ne0$ and $b\ne0$. Since $\theta^3=-a\theta-b$, we have $\theta f'(\theta)=3\theta^3+a\theta=
-2a\theta-3b$, hence $N(\theta)\,N(f'(\theta))=N(-2a\theta-3b)$. Here $N(\theta)=\theta_1\theta_2\theta_3=-b$, and, with
$c=-\frac{3b}{2a}$ and $\prod_i(c-\theta_i)=f(c)$,
\begin{align*}
N(-2a\theta-3b)&=\prod_i\bigl(-2a(\theta_i-c)\bigr)\\&=(-2a)^3(-1)^3f(c)=8a^3f(c)\\
&=8a^3\Bigl(-\frac{27b^3}{8a^3}-\frac{3b}2+b\Bigr)\\&=-27b^3-4a^3b.
\end{align*}
So $N(f'(\theta))=\frac{-b(27b^2+4a^3)}{-b}=4a^3+27b^2$, and Proposition~\ref{prop:disc}(c) with $n=3$ gives
$\operatorname{disc}(f)=(-1)^3N(f'(\theta))=-4a^3-27b^2$. Both sides are polynomials in $(a,b)$ (the left one is a
symmetric polynomial in the roots, hence a polynomial in the coefficients), equal on the dense set $ab\ne0$, so they are
equal everywhere. Check: $X^3-X$ has roots $0,1,-1$ and discriminant $1\cdot1\cdot4=4=-4(-1)^3$.

For $f=X^3-X-1$: $-4(-1)^3-27=-23$, which is squarefree, so $\mathcal O_K=\Z[\theta]$ and $d_K=-23$ by
Proposition~\ref{prop:index}. Since $d_K<0$, $r_2$ is odd (Proposition~\ref{prop:sign}); with $r_1+2r_2=3$ this
forces $r_2=1$, $r_1=1$: one real embedding.
\end{solution}

\begin{solution}{ex:normform}
\idea The norm is the determinant of multiplication; and a unit of norm $1$ gives, by its powers, infinitely many
elements of norm $1$, which are infinitely many solutions.
\writeup With $\theta^3=2$: $\alpha\cdot1=a+b\theta+c\theta^2$, $\alpha\theta=2c+a\theta+b\theta^2$,
$\alpha\theta^2=2b+2c\theta+a\theta^2$. The matrix with these columns is the one displayed in the text, and expanding
the determinant gives $a(a^2-2bc)-2c(ab-2c^2)+2b(b^2-ac)=a^3+2b^3+4c^3-6abc$. The unit $u=\theta-1$ has norm
$(-1)^3+2\cdot1^3-0=1$, so $N(u^k)=1$ for all $k\in\Z$, and the $u^k$ are pairwise distinct (as $|u|=0.26\ldots\ne1$),
giving infinitely many solutions. For $k=3$: $(\theta-1)^3=\theta^3-3\theta^2+3\theta-1=1+3\theta-3\theta^2$, so
$(a,b,c)=(1,3,-3)$, and indeed $1+54-108+54=1$.
\end{solution}

\begin{solution}{ex:discbinom}
\idea Proposition~\ref{prop:disc}(c) gives the general formula; for $X^n-a$ the derivative is a monomial, whose norm is
easy.
\writeup By Proposition~\ref{prop:disc}(c), $\operatorname{disc}(f)=\prod_{i<j}(\theta_i-\theta_j)^2=(-1)^{n(n-1)/2}
\prod_if'(\theta_i)$, and this computation only used that $f$ is monic with distinct roots. For $f=X^n-a$:
$f'(\theta_i)=n\theta_i^{n-1}$, and $\prod_i\theta_i=(-1)^nf(0)=(-1)^{n+1}a$. So
$\prod_if'(\theta_i)=n^n\bigl((-1)^{n+1}a\bigr)^{n-1}=n^n(-1)^{(n+1)(n-1)}a^{n-1}$, and
$\operatorname{disc}=(-1)^{n(n-1)/2+(n^2-1)}n^na^{n-1}$. The exponent $\frac{n(n-1)}2+n^2-1$ has the same parity as
$\frac{(n-1)(n-2)}2$, since their difference is $\frac{n^2-n-(n^2-3n+2)}2+n^2-1=n-1+n^2-1=n^2+n-2$, which is even. For
$X^3-2$: $(-1)^1\cdot27\cdot4=-108$, as in Example~\ref{exa:cuberoot2}.
\end{solution}

\begin{solution}{ex:kronecker}
\idea Finitely many integer polynomials of degree $n$ have bounded coefficients; the powers of $\alpha$ all have such
characteristic polynomials, so two powers share one, and a pigeonhole argument closes the loop.
\writeup Let $K=\Q(\alpha)$, $n=[K:\Q]$. For $k\ge1$, $\chi_{\alpha^k}=\prod_i(X-\sigma_i(\alpha)^k)$ has integer
coefficients (Proposition~\ref{prop:minpolyZ}), and its coefficient of $X^{n-j}$ is bounded by $\binom nj$ in absolute
value since all roots have modulus $\le1$. So the $\chi_{\alpha^k}$ take finitely many values, and so do the sets of
their roots; in particular the powers $\alpha^k$ (roots of these polynomials) take finitely many values. So
$\alpha^k=\alpha^m$ for some $k<m$, and if $\alpha\ne0$ then $\alpha^{m-k}=1$. Integrality is needed: $\alpha=
\frac{3+4i}5$ has conjugates $\frac{3\pm4i}5$ of modulus $1$, but it is not a root of unity (its powers have growing
denominators $5^k$).
\end{solution}

\begin{solution}{ex:sqrt23}
\idea Eliminate the square roots by squaring; the degree count shows that $\alpha$ generates the whole biquadratic field.
\writeup $\alpha^2=5+2\sqrt6$, $(\alpha^2-5)^2=24$, so $\alpha^4-10\alpha^2+1=0$. As $\Q(\alpha)\subseteq\Q(\sqrt2,\sqrt3)$,
which has degree $4$ (Tome~2, Chapter~8: $\sqrt3\notin\Q(\sqrt2)$), and $\alpha^3=11\sqrt2+9\sqrt3$ gives
$\sqrt2=\frac{\alpha^3-9\alpha}2$ and $\sqrt3=\frac{11\alpha-\alpha^3}2$, the two fields are equal, of degree $4$, and
$X^4-10X^2+1$ is the minimal polynomial. The element $\gamma=\frac{\sqrt2+\sqrt6}2$ has $\gamma^2=\frac{8+2\sqrt{12}}4=
2+\sqrt3$, an algebraic integer, so $\gamma$ is integral (Exercise~\ref{ex:roots}). It is not in
$\Z[\sqrt2,\sqrt3]=\Z+\Z\sqrt2+\Z\sqrt3+\Z\sqrt6$, because its coordinates in the basis $1,\sqrt2,\sqrt3,\sqrt6$ are
$0,\frac12,0,\frac12$.
\end{solution}

\begin{solution}{ex:biquadratic}
\idea The field has three quadratic subfields, and the trace from $K$ to each of them (the sum over the two
automorphisms fixing it) sends integers to integers of the subfield. These three conditions pin down the
denominators. Then the discriminant is a $4\times4$ determinant of traces.
\writeup The automorphisms of $K$ change the signs of $\sqrt2$ and $\sqrt3$ independently. Let
$x=a+b\sqrt2+c\sqrt3+d\sqrt6\in\mathcal O_K$ with $a,b,c,d\in\Q$. Its trace to $\Q(\sqrt2)$ (sum of $x$ and its image
under $\sqrt3\mapsto-\sqrt3$) is $2a+2b\sqrt2$, an integer of $\Q(\sqrt2)$, so $2a,2b\in\Z$; to $\Q(\sqrt3)$ it is
$2a+2c\sqrt3$, so $2c\in\Z$; to $\Q(\sqrt6)$ it is $2a+2d\sqrt6$, so $2d\in\Z$. Write $a,b,c,d=\frac{A}2,\frac B2,\frac C2,
\frac D2$. Now use the norm to $\Q(\sqrt2)$, $x\cdot x'$ with $x'=a+b\sqrt2-c\sqrt3-d\sqrt6$: it equals
$(a+b\sqrt2)^2-3(c+d\sqrt2)^2=a^2+2b^2-3c^2-6d^2+(2ab-6cd)\sqrt2$, which must be in $\Z[\sqrt2]$:
$A^2+2B^2-3C^2-6D^2\equiv0\pmod4$ and $AB-3CD\equiv0\pmod2$. The same with $\Q(\sqrt3)$ (norm
$(a+c\sqrt3)^2-2(b+d\sqrt3)^2$) gives $A^2-2B^2+3C^2-6D^2\equiv0\pmod4$ and $AC-2BD\equiv0\pmod2$. Subtracting the
two first congruences, $4B^2-6C^2\equiv0\pmod4$, so $C$ is even. Then $A^2+2B^2-6D^2\equiv0\pmod4$ forces $A$ even, and
then $2B^2\equiv6D^2\pmod4$, i.e.\ $B\equiv D\pmod2$. So $x\in\Z+\Z\sqrt2+\Z\sqrt3+\Z\frac{\sqrt2+\sqrt6}2$, and
conversely all four basis elements are integral (Exercise~\ref{ex:sqrt23}). For the discriminant,
$\operatorname{disc}(1,\sqrt2,\sqrt3,\sqrt6)=\det(\operatorname{Tr}(\omega_i\omega_j))=\det\operatorname{diag}(4,8,12,24)=
9216$, and the basis $1,\sqrt2,\sqrt3,\frac{\sqrt2+\sqrt6}2$ is obtained by a matrix of determinant $\frac12$, so
$d_K=\frac{9216}4=2304=2^8\cdot3^2$. Finally $\operatorname{disc}(X^4-10X^2+1)=147456=2304\cdot64$ (from the roots
$\pm\sqrt2\pm\sqrt3$, or with the notebook), so the index of $\Z[\sqrt2+\sqrt3]$ is $8$.
\end{solution}

\begin{solution}{ex:cuberoot10}
\idea The index formula leaves only the primes $2,3,5$, since $\operatorname{disc}(X^3-10)=-27\cdot100=-2700=-2^23^35^2$.
Eisenstein removes $2$ and $5$; for $3$, the element $\beta$ shows the index is at least $3$, and a trace argument shows
that $\Z[\theta,\beta]$ is already $3$-maximal.
\writeup By Exercise~\ref{ex:discbinom}, $\operatorname{disc}(X^3-10)=-27\cdot100=-2700$, so $[\mathcal O_K:\Z[\theta]]^2$
divides $2^2\cdot3^3\cdot5^2$, and the index divides $2\cdot3\cdot5$. $X^3-10$ is $2$- and $5$-Eisenstein, so the
index divides $3$ (a square dividing $3^3$ divides $3^2$); it is $3$, since $\beta\in\mathcal O_K\setminus\Z[\theta]$
(Example~\ref{exa:cuberoot10}). Let $M=\Z+\Z\theta+\Z\beta$. Then $\Z[\theta]\subseteq M\subseteq\mathcal O_K$, since
$\theta^2=3\beta-1-\theta\in M$, and the change of basis from $1,\theta,\theta^2$ to $1,\theta,\beta$ has determinant
$\frac13$, so $[M:\Z[\theta]]=3$. Therefore $[\mathcal O_K:M]=[\mathcal O_K:\Z[\theta]]/[M:\Z[\theta]]=1$:
$\mathcal O_K=M$, an integral basis is $1,\theta,\beta$, and $d_K=\operatorname{disc}(1,\theta,\theta^2)/9=-300$
(Proposition~\ref{prop:disc}(b)). For the record, $M$ is visibly a ring: $\theta\beta=\frac{\theta+\theta^2+10}3=\beta+3$
and, from $(1+\theta+\theta^2)^2=1+2\theta+3\theta^2+2\theta^3+\theta^4=21+12\theta+3\theta^2$,
$\beta^2=\frac{7+4\theta+\theta^2}3=\beta+2+\theta$.
\end{solution}

\begin{solution}{ex:totallyreal}
\idea The triple-angle formula gives the equation. The discriminant $81>0$ gives the signature, Eisenstein at $3$ after a
translation gives the ring of integers, and the other roots come from $\cos$ of the other angles $\frac{4\pi}9,\frac{8\pi}9$.
\writeup With $t=\frac{2\pi}9$ and $\theta=2\cos t$: $\theta^3-3\theta=8\cos^3t-6\cos t=2\cos3t=2\cos\frac{2\pi}3=-1$. The
polynomial $X^3-3X+1$ has no rational (hence integer) root ($\pm1$ fail), so it is irreducible, and its discriminant is
$-4(-3)^3-27=81$ (Exercise~\ref{ex:disccubic}). As $81>0$, $r_2$ is even, so $r_2=0$: the field is totally real. The
index divides $9$. The translate $\theta+1$ has minimal polynomial $(X-1)^3-3(X-1)+1=X^3-3X^2+3$, which is
$3$-Eisenstein, so $3$ does not divide the index: $\mathcal O_K=\Z[\theta]$, $d_K=81$. Finally $\theta^2-2=4\cos^2t-2=
2\cos2t$ and $2-\theta-\theta^2=-2\cos t-2\cos2t=2\cos4t$ (since $\cos t+\cos2t+\cos4t=0$: the sum of the three roots is
$0$, and $2\cos4t=2\cos\frac{8\pi}9$ is the third root of the same equation, by the triple-angle formula with $3\cdot
\frac{8\pi}9=\frac{8\pi}3\equiv\frac{2\pi}3$). All three roots lie in $\Q(\theta)$, so it is the splitting field: Galois,
with cyclic group of order $3$.
\end{solution}

\begin{solution}{ex:dedekind}
\idea Show integrality of $\beta$ by finding its minimal polynomial or, better, by showing $\Z+\Z\theta+\Z\beta$ is a ring
(Proposition~\ref{prop:intcrit}(c)). Then the index formula, with $503$ prime, finishes.
\writeup $f$ has no root among $\pm1,\pm2,\pm4,\pm8$, so it is irreducible. Its discriminant can be computed from the
general cubic formula $\operatorname{disc}(X^3+bX^2+cX+d)=b^2c^2-4c^3-4b^3d-27d^2+18bcd$ with $(b,c,d)=(-1,-2,-8)$:
$4+32-32-1728-288=-2012$. Let $M=\Z+\Z\theta+\Z\beta$. Using $\theta^3=\theta^2+2\theta+8$:
\[
\theta\beta=\frac{\theta^3+\theta^2}2=\theta^2+\theta+4=2\beta+4,\qquad
\beta^2=\frac{\theta^4+2\theta^3+\theta^2}4.
\]
Now $\theta^4=\theta^3+2\theta^2+8\theta=3\theta^2+10\theta+8$, so $\theta^4+2\theta^3+\theta^2=3\theta^2+10\theta+8+
2\theta^2+4\theta+16+\theta^2=6\theta^2+14\theta+24$, and $\beta^2=\frac{6\theta^2+14\theta+24}4=3\beta+2\theta+6$.
So $M$ is stable under multiplication by $\theta$ (using $\theta^2=2\beta-\theta$) and by $\beta$: it is a ring,
finitely generated as a group, so its elements are algebraic integers by Proposition~\ref{prop:intcrit}(c). The index
$[M:\Z[\theta]]$ is $2$ (the basis change $1,\theta,\theta^2\to1,\theta,\beta$ has determinant $\frac12$), so
$\operatorname{disc}(1,\theta,\beta)=-2012/4=-503$. Since $503$ is prime (it has no prime factor up to $22$), the index
formula gives $[\mathcal O_K:M]=1$: $\mathcal O_K=M$ and $d_K=-503$. An exercise of Chapter~2 shows that no single
$\gamma$ generates $\mathcal O_K$; the reason is that $2$ splits into three different prime ideals of norm $2$,
while $\mathbb F_2$ has only two elements to be the roots of $m_\gamma\bmod2$.
\end{solution}

\begin{solution}{ex:traceform}
\idea Diagonalise the trace form using the embeddings as coordinates: real embeddings give squares, and each complex pair
gives a hyperbolic plane $u^2-v^2$.
\writeup Let $\sigma_1,\dots,\sigma_{r_1}$ be the real embeddings and $\tau_1,\bar\tau_1,\dots,\tau_{r_2},\bar\tau_{r_2}$
the complex ones. The map
\[
\Phi:x\mapsto\bigl(\sigma_1(x),\dots,\sigma_{r_1}(x),\operatorname{Re}\tau_1(x),\operatorname{Im}\tau_1(x),\dots\bigr)
\]
from $K$ to $\R^n$ is $\Q$-linear, and its real extension $K\otimes\R\to\R^n$ is a bijection: its
matrix in a basis $\omega_j$ has determinant $\pm(2i)^{-r_2}\det(\sigma_i(\omega_j))\ne0$ (replace the two rows $\tau,\bar
\tau$ by $\frac{\tau+\bar\tau}2,\frac{\tau-\bar\tau}{2i}$). In these coordinates,
$\operatorname{Tr}(x^2)=\sum_k\sigma_k(x)^2+\sum_l\bigl(\tau_l(x)^2+\overline{\tau_l(x)}^2\bigr)=\sum_ky_k^2+
\sum_l2(u_l^2-v_l^2)$, writing $\tau_l(x)=u_l+iv_l$. This form has $r_1+r_2$ positive and $r_2$ negative squares. The
determinant of a real quadratic form has the sign $(-1)^{\text{number of negative squares}}$, and in the basis
$(\omega_j)$ the matrix of $\operatorname{Tr}(x^2)$ is $(\operatorname{Tr}(\omega_i\omega_j))$, with determinant $d_K$:
so $d_K$ has sign $(-1)^{r_2}$. Finally the form is positive definite on $K\otimes\R$ iff $r_2=0$; and positivity on
$K^\times$ is equivalent to positivity on $K\otimes\R$ (a negative value on a real vector is approached by rational
vectors, by continuity), which gives the last claim.
\end{solution}

\section*{Sources and further reading}

\begin{itemize}
\item J.~Neukirch, \emph{Algebraic Number Theory}, Grundlehren~322, Springer, 1999, Chapter~I, \S\S1--2. Integrality and
discriminants in the language of rings.
\item D.~A.~Marcus, \emph{Number Fields}, 2nd ed., Universitext, Springer, 2018, Chapter~2. The source for the Eisenstein
argument of Proposition~\ref{prop:eisenstein}, with many exercises on cubic and cyclotomic fields.
\item P.~Samuel, \emph{Algebraic Theory of Numbers}, Dover, 2008, Chapter~2. Short and clear.
\item J.~S.~Milne, \emph{Algebraic Number Theory}, course notes (free online), Chapter~2.
\item L.~Stickelberger, ``\"Uber eine neue Eigenschaft der Diskriminanten algebraischer Zahlk\"orper'', \emph{Verhandlungen
des ersten Internationalen Mathematiker-Kongresses}, Z\"urich, 1897. The proof of Theorem~\ref{thm:stickelberger} given
here is due to Schur (1929).
\item R.~Dedekind, ``\"Uber den Zusammenhang zwischen der Theorie der Ideale und der Theorie der h\"oheren Kongruenzen'',
\emph{Abh.\ K\"onigl.\ Ges.\ Wiss.\ G\"ottingen} 23 (1878). The cubic field of Exercise~\ref{ex:dedekind}.
\item H.~Cohen, \emph{A Course in Computational Algebraic Number Theory}, GTM~138, Springer, 1993, Chapter~6. How
computers find $\mathcal O_K$ (the Round~2 and Round~4 algorithms), used by PARI/GP in the notebook.
\end{itemize}
